\begin{figure}[htbp]
    \centering
    \subfloat[wenige kurze Testcases: $T$=10, $I$=100]{%
        \resizebox{0.45\textwidth}{!}{%
        \benchmarkplotOverhead{
            (0, 122.0)
            (20, 122.4)
            (40, 124.5)
            (60, 129.2)
            (80, 120.6)
            (100, 121.6)}
            {
                (0, 3.6)
                (20, 3.6)
                (40, 3.6)
                (60, 3.6)
                (80, 3.5)
                (100, 3.6)
            }{
                (0, 210.2)
                (20, 138.8)
                (40, 138.9)
                (60, 134.9)
                (80, 134.9)
                (100, 9.6)
            }{
                (0, 39.6)
                (20, 41.1)
                (40, 41.6)
                (60, 41.5)
                (80, 41.1)
                (100, 39.3)
            }%
        }%
    }
    \hfill
    \subfloat[viele kurze Testcases: $T$=100, $I$=10]{%
        \resizebox{0.45\textwidth}{!}{%
        \benchmarkplotOverhead{
            (0, 161.8)
            (20, 165.9)
            (40, 162.5)
            (60, 163.0)
            (80, 161.2)
            (100, 164.2)
        }{
            (0, 382.8)
            (20, 370.0)
            (40, 352.5)
            (60, 336.6)
            (80, 367.9)
            (100, 177.9)
        }{}{}%
        }%
    }
    \vspace{-0.5em}
    \subfloat[wenige mittellange Testcases: $T$=10, $I$=1.000.000]{%
        \resizebox{0.45\textwidth}{!}{%
        \benchmarkplotOverhead{
            (0, 456.6)
            (20, 427.6)
            (40, 372.0)
            (60, 443.3)
            (80, 357.4)
            (100, 363.5)
        }{
            (0, 553.0)
            (20, 523.3)
            (40, 472.8)
            (60, 400.1)
            (80, 347.3)
            (100, 175.0)
        }{}{}%
        }%
    }
    \hfill
    \subfloat[viele mittellange Testcases: $T$=1.000, $I$=10.000]{%
        \resizebox{0.45\textwidth}{!}{%
        \benchmarkplotOverhead{
            (0, 783.7)
            (20, 789.6)
            (40, 778.1)
            (60, 843.8)
            (80, 803.6)
            (100, 787.3)
        }{
            (0, 1599.1)
            (20, 1318.1)
            (40, 1063.7)
            (60, 830.1)
            (80, 596.4)
            (100, 215.2)
        }{}{}%
        }%
    }
    \vspace{-0.5em}
    \subfloat[wenige langsame Testcases: $T$=10, $I$=10.000.000]{%
        \resizebox{0.45\textwidth}{!}{%
        \benchmarkplotOverhead{
            (0, 2580.0)
            (20, 2418.9)
            (40, 2739.4)
            (60, 2611.9)
            (80, 2666.4)
            (100, 2378.4)
        }{
            (0, 3070.7)
            (20, 2237.5)
            (40, 1755.5)
            (60, 1239.1)
            (80, 783.0)
            (100, 228.3)
        }{}{}%
        }%
    }
    \hfill
    \subfloat[viele langsame Testcases: $T$=10.000, $I$=10.000]{%
        \resizebox{0.45\textwidth}{!}{%
        \benchmarkplotOverhead{
            (0, 7118.8)
            (20, 7074.7)
            (40, 7068.8)
            (60, 7080.6)
            (80, 7115.5)
            (100, 7116.6)
        }{
            (0, 21777.5)
            (20, 17293.3)
            (40, 13253.1)
            (60, 9229.9)
            (80, 5158.9)
            (100, 831.9)
        }{}{}%
        }%
    }
    \caption{Detailansicht zu den Laufzeiten einzelner Komponenten des Verfahrens. Sämtliche dieser Übersicht zugrunde liegenden Daten können in dem Repository der Benchmark (siehe \chapref{Links}) eingesehen und überprüft werden.}
    \label{fig:benchmark_overhead}
\end{figure}
